
%%%%%%%%%%%%%%%%%%%%%%%%%%%%%%%%%%%%%%
%%            Introdução            %%
%%%%%%%%%%%%%%%%%%%%%%%%%%%%%%%%%%%%%%

\section{Introdução}
\label{sec:introducao}

\begin{frame}{Introdução}

	\begin{itemize}
		\item \textbf{Tomada de decisão:} é um processo de escolha entre duas ou mais alternativas, que envolve a avaliação de riscos e benefícios, com o objetivo de maximizar os benefícios e minimizar os riscos.
		\item \textbf{Análise de Dados:} é o processo de coleta, organização, análise e interpretação de dados, com o objetivo de extrair informações úteis.
	\end{itemize}

	\bigskip

	\setlength{\leftskip}{3cm}

	\quad Os processos de tomada de decisão e análise de dados são fundamentais para o sucesso de qualquer organização, seja ela uma empresa, uma instituição de ensino ou uma instituição governamental.

	\medskip

	\quad Para realizar a tarefa de análise de dados, os profissionais e pesquisadores utilizam ferramentas computacionais, que permitem a manipulação de dados, a geração de gráficos e a realização de cálculos estatísticos.

	\setlength{\leftskip}{0pt}


\end{frame}

\subsection{Objetivos}
\label{subsec:objetivos}

\begin{frame}{Objetivos}

	\quad O objetivo deste trabalho é o estudo e desenvolvimento de um projeto de software com a finalidade de análise e visualização de dados, que seja modular e facilmente extensível, para atender a vários propósitos e fontes de dados.

\end{frame}

\subsubsection{Objetivo Geral}
\label{subsubsec:objetivo-geral}

\begin{frame}{Objetivo Geral}

	\textbf{Objetivo Geral:} Desenvolver um sistema de análise de dados para a tomada de decisão, utilizando a linguagem de programação Python.

\end{frame}

\subsubsection{Objetivos Específicos}
\label{subsubsec:objetivos-especificos}

\begin{frame}{Objetivos Específicos}

	\begin{itemize}
		\item Identificar diferentes problemas de análise de dados, tipos e fontes de dados.
		\item Modelar um projeto de software que atenda as necessidades identificadas, utilizando os conceitos de \textbf{ES}.
		\item Implementar o software modelado.
		\item Testar e validar o software aplicando-o em algumas bases de dados.
		\item Documentar o software para futuras modificações.
    \end{itemize}

\end{frame}

\subsection{Justificativa}
\label{subsec:justificativa}

\begin{frame}{Justificativa}

    A seguir são apresentadas algumas justificativas para a escolha do tema deste trabalho.

    \medskip

    \begin{itemize}
        \item  A tarefa de análise de dados é complexa e exige o conhecimento e emprego de diversas técnicas de processamento.
        \item O processo exploratório da análise de dados requer o emprego de técnicas de visualização de dados, que permitem a compreensão dos dados e a identificação de padrões e tendências.
        \item Frequentemente, a escolha das técnicas é baseada em experiência e intuição, necessitando de uma série de testes e ajustes dos parâmetros para a obtenção de resultados satisfatórios.
        \item A aplicação de determinadas técnica são realizadas em conjunto, e a escolha da ordem de aplicação das técnicas é um desafio, pois a ordem de aplicação das técnicas pode influenciar no resultado.
    \end{itemize}

\end{frame}
